% Sounio: Epistemic Types for Scientific Computing
% ArXiv Preprint - Computer Science > Programming Languages
%
% Target: cs.PL, cs.SE, stat.CO
% Keywords: uncertainty quantification, type systems, scientific computing, GUM

\documentclass[11pt,a4paper]{article}

% Packages
\usepackage[utf8]{inputenc}
\usepackage[T1]{fontenc}
\usepackage{amsmath,amssymb,amsthm}
\usepackage{algorithm}
\usepackage{algpseudocode}
\usepackage{booktabs}
\usepackage{graphicx}
\usepackage{hyperref}
\usepackage{listings}
\usepackage{xcolor}
\usepackage{natbib}
\usepackage{doi}

% Code listing style
\definecolor{codebg}{RGB}{248,248,248}
\definecolor{keyword}{RGB}{0,0,180}
\definecolor{comment}{RGB}{0,128,0}
\definecolor{string}{RGB}{163,21,21}

\lstdefinelanguage{sounio}{
  keywords={fn, let, var, if, else, while, for, return, struct, enum, match, with, kernel, linear, use},
  keywordstyle=\color{keyword}\bfseries,
  ndkeywords={Knowledge, f64, i32, bool, string, IO, Alloc, GPU},
  ndkeywordstyle=\color{blue},
  sensitive=true,
  comment=[l]{//},
  commentstyle=\color{comment}\itshape,
  stringstyle=\color{string},
  morestring=[b]",
}

\lstset{
  language=sounio,
  basicstyle=\ttfamily\small,
  backgroundcolor=\color{codebg},
  frame=single,
  numbers=left,
  numberstyle=\tiny\color{gray},
  breaklines=true,
  captionpos=b,
}

% Theorems
\newtheorem{theorem}{Theorem}
\newtheorem{lemma}[theorem]{Lemma}
\newtheorem{definition}{Definition}
\newtheorem{proposition}{Proposition}

% Title
\title{Sounio: Epistemic Types for Scientific Computing\\with Native Uncertainty Quantification}

\author{
  Demetrios Chiuratto Agourakis\\
  \texttt{demetrios@sounio-lang.org}\\
  \url{https://souniolang.org}
}

\date{January 2026}

\begin{document}

\maketitle

\begin{abstract}
We present Sounio, a systems programming language with native support for
\emph{epistemic types}---type-level representations of uncertain values that
automatically propagate measurement uncertainty through computations.
The \texttt{Knowledge<T>} type encapsulates a value, its variance, confidence
distribution, and provenance chain, enabling GUM-compliant (Guide to the
Expression of Uncertainty in Measurement) uncertainty propagation without
library overhead. We describe the type system, effect-based safety guarantees,
and compile-time optimizations that achieve 2$\times$ speedup over Julia's
\texttt{Measurements.jl} and 30$\times$ over Python's \texttt{uncertainties}
package. We demonstrate applications in pharmacokinetics, climate modeling,
and neuroimaging, where correct uncertainty handling is essential for
regulatory compliance and scientific reproducibility.
\end{abstract}

\section{Introduction}
\label{sec:intro}

Scientific computing frequently involves uncertain quantities: measured values
with instrument precision, model parameters with confidence intervals, and
predictions with error bounds. Yet mainstream programming languages treat
numerical values as exact, forcing developers to either:

\begin{enumerate}
  \item Ignore uncertainty (compromising scientific validity)
  \item Track uncertainty manually (error-prone, tedious)
  \item Use library wrappers (performance overhead, API friction)
\end{enumerate}

The Guide to the Expression of Uncertainty in Measurement (GUM)
\citep{jcgm2008gum} provides a rigorous framework for uncertainty propagation,
but implementing GUM in general-purpose languages requires discipline that
compilers cannot enforce.

We introduce \textbf{Sounio}, a systems programming language where uncertainty
is a \emph{first-class type-level concept}. The \texttt{Knowledge<T>} type
carries not just a value but also:

\begin{itemize}
  \item \textbf{Variance}: Uncertainty in the GUM sense ($u^2$)
  \item \textbf{Confidence}: A Beta distribution over credibility
  \item \textbf{Provenance}: Audit trail for regulatory compliance
\end{itemize}

Arithmetic on \texttt{Knowledge<T>} automatically propagates uncertainty using
first-order Taylor series (GUM linear method) or Monte Carlo sampling (GUM
Supplement 1), with the choice configurable per-computation.

\paragraph{Contributions}
\begin{itemize}
  \item A type system for epistemic values with compile-time effect tracking
        (Section~\ref{sec:type-system})
  \item GUM-compliant uncertainty propagation integrated into the compiler's
        intermediate representation (Section~\ref{sec:propagation})
  \item Benchmarks showing 2--30$\times$ speedup over existing solutions
        (Section~\ref{sec:evaluation})
  \item Case studies in pharmacokinetics, climate science, and neuroimaging
        (Section~\ref{sec:applications})
\end{itemize}

\section{Background}
\label{sec:background}

\subsection{The GUM Framework}

The GUM \citep{jcgm2008gum} defines standard uncertainty $u(x)$ as the
positive square root of variance:
\[
u(x) = \sqrt{\text{Var}(x)}
\]

For a function $y = f(x_1, \ldots, x_N)$ of uncorrelated inputs, the
\emph{combined standard uncertainty} is:
\begin{equation}
  u_c^2(y) = \sum_{i=1}^{N} \left(\frac{\partial f}{\partial x_i}\right)^2 u^2(x_i)
  \label{eq:gum-propagation}
\end{equation}

This first-order Taylor expansion (the ``law of propagation of uncertainty'')
is exact for linear functions and approximate otherwise.

\subsection{Existing Approaches}

\paragraph{Python: uncertainties}
The \texttt{uncertainties} package \citep{lebigot2024uncertainties} provides
\texttt{ufloat} objects with automatic differentiation for Equation~\ref{eq:gum-propagation}.
However, operator overloading in Python incurs significant runtime overhead
(see Section~\ref{sec:evaluation}).

\paragraph{Julia: Measurements.jl}
Julia's \texttt{Measurements.jl} \citep{giordano2016measurements} uses dual
numbers for uncertainty tracking, achieving better performance through JIT
compilation. However, it remains a library with no compile-time guarantees.

\paragraph{Domain-Specific Languages}
Languages like Stan \citep{carpenter2017stan} and BUGS model uncertainty
probabilistically but are not general-purpose systems languages.

\subsection{Limitations of Library Approaches}

Library-based uncertainty tracking has three fundamental limitations:

\begin{enumerate}
  \item \textbf{No static guarantees}: A function accepting \texttt{float64}
        can silently drop uncertainty.
  \item \textbf{Runtime overhead}: Wrapper types require heap allocation,
        indirection, and runtime dispatch.
  \item \textbf{No provenance}: Audit trails for regulatory compliance must
        be implemented separately.
\end{enumerate}

Sounio addresses all three by making uncertainty a compile-time type property.

\section{The Sounio Type System}
\label{sec:type-system}

\subsection{The Knowledge Type}

The core epistemic type is:

\begin{lstlisting}[caption={Knowledge type definition}]
struct Knowledge<T> {
    value: T,                    // Point estimate
    variance: f64,               // Uncertainty (sigma^2)
    confidence: BetaConfidence,  // Credibility distribution
    provenance: Provenance,      // Audit trail
}
\end{lstlisting}

\texttt{BetaConfidence} represents epistemic confidence as a Beta distribution,
supporting Bayesian updating:

\begin{lstlisting}[caption={Confidence representation}]
struct BetaConfidence {
    alpha: f64,  // Successes + prior
    beta: f64,   // Failures + prior
}

impl BetaConfidence {
    fn mean(&self) -> f64 { self.alpha / (self.alpha + self.beta) }
    fn update(&self, successes: i64, failures: i64) -> BetaConfidence { ... }
}
\end{lstlisting}

\subsection{Type-Level Effect Tracking}

Sounio uses an algebraic effect system to track computational side effects.
The effect \texttt{Epistemic} marks functions that introduce or propagate
uncertainty:

\begin{lstlisting}[caption={Effect-tracked measurement function}]
fn measure(instrument: &Instrument) -> Knowledge<f64> with IO, Epistemic {
    let raw = instrument.read()  // IO effect
    Knowledge::new(
        value: raw,
        variance: instrument.precision_squared(),
        confidence: BetaConfidence::from_calibration(instrument),
        provenance: Provenance::measurement(instrument)
    )
}
\end{lstlisting}

A function without \texttt{with Epistemic} cannot silently discard uncertainty:

\begin{lstlisting}[caption={Type error: uncertainty lost}]
fn compute_exact(x: f64) -> f64 { x * 2.0 }

fn main() with IO {
    let k = Knowledge::measured(5.0, 0.01, "scale")
    let bad = compute_exact(*k.get())  // Warning: uncertainty discarded
    let good = k.scale(2.0)            // OK: uncertainty preserved
}
\end{lstlisting}

\subsection{Subtyping and Coercion}

\texttt{Knowledge<T>} is not a subtype of \texttt{T}. Explicit extraction is
required, generating compiler warnings when uncertainty is discarded:

\begin{definition}[Epistemic Extraction]
The operation $\text{extract}: \texttt{Knowledge<T>} \to \texttt{T}$ is
\emph{effect-polymorphic}: it is always permitted but annotates the call
site with a ``discarded uncertainty'' warning unless suppressed by
\texttt{\#[allow(uncertainty\_loss)]}.
\end{definition}

\subsection{Provenance Tracking}

Every \texttt{Knowledge} value carries a \texttt{Provenance} chain:

\begin{lstlisting}[caption={Provenance structure}]
enum Source {
    Measurement { instrument: String, timestamp: i64 },
    Computed { operation: String },
    Assertion { author: String },
    External { url: String, doi: Option<String> },
}

struct Provenance {
    sources: Vec<Source>,
    merkle_hash: [u8; 32],  // Tamper-evident hash
}
\end{lstlisting}

This supports regulatory requirements (FDA 21 CFR Part 11, ISO 17025) for
audit trails in validated computational systems.

\section{Uncertainty Propagation}
\label{sec:propagation}

\subsection{First-Order (Linear) Propagation}

For arithmetic operations, Sounio implements Equation~\ref{eq:gum-propagation}
at the IR level. Consider multiplication:

\begin{align}
  z &= x \cdot y \\
  u^2(z) &= y^2 u^2(x) + x^2 u^2(y) \quad \text{(assuming independence)}
\end{align}

The compiler transforms:

\begin{lstlisting}
let z = x * y  // where x, y: Knowledge<f64>
\end{lstlisting}

into:

\begin{lstlisting}
let z_val = x.value * y.value
let z_var = y.value^2 * x.variance + x.value^2 * y.variance
let z = Knowledge { value: z_val, variance: z_var, ... }
\end{lstlisting}

\subsection{Transcendental Functions}

For $y = f(x)$, the delta method gives:
\[
  u^2(y) \approx \left(\frac{df}{dx}\right)^2 u^2(x)
\]

Sounio provides built-in derivatives for standard functions:

\begin{center}
\begin{tabular}{lll}
\toprule
Function & $f(x)$ & $u^2(f(x))$ \\
\midrule
\texttt{sqrt} & $\sqrt{x}$ & $\frac{u^2(x)}{4x}$ \\
\texttt{exp} & $e^x$ & $e^{2x} u^2(x)$ \\
\texttt{ln} & $\ln x$ & $\frac{u^2(x)}{x^2}$ \\
\texttt{sin} & $\sin x$ & $\cos^2(x) u^2(x)$ \\
\texttt{cos} & $\cos x$ & $\sin^2(x) u^2(x)$ \\
\bottomrule
\end{tabular}
\end{center}

\subsection{Confidence Combination}

When combining \texttt{Knowledge} values, confidence is updated using:

\begin{proposition}[Confidence Product Rule]
For independent measurements $x$ and $y$ with Beta-distributed confidences,
the confidence of $z = f(x, y)$ is:
\[
  \text{BetaConfidence}(z) = \text{combine}(\text{conf}(x), \text{conf}(y))
\]
where $\text{combine}$ pools evidence under the assumption of conditional
independence.
\end{proposition}

\subsection{Scientific IR (SIR) Optimizations}

Sounio's compiler includes a Scientific IR pass that optimizes epistemic
computations:

\begin{enumerate}
  \item \textbf{Variance fusion}: $u^2(x + y + z)$ computed in one pass
  \item \textbf{Exact propagation}: Operations on \texttt{Knowledge::exact}
        skip variance computation
  \item \textbf{SIMD vectorization}: Value and variance computed in parallel
  \item \textbf{Provenance elision}: When not required, provenance chains
        are not allocated
\end{enumerate}

\section{Evaluation}
\label{sec:evaluation}

\subsection{Benchmark Suite}

We compare Sounio against Python's \texttt{uncertainties} and Julia's
\texttt{Measurements.jl} on five benchmarks:

\begin{enumerate}
  \item \textbf{GUM Chain}: 100K arithmetic operations with uncertainty
  \item \textbf{Monte Carlo}: 100K portfolio simulations
  \item \textbf{Cockcroft-Gault}: 10K pharmacokinetic calculations
  \item \textbf{Matrix}: 50$\times$50 matrix-vector multiply with uncertainty
  \item \textbf{Transcendental}: 10K chained sin/cos/exp/log/sqrt
\end{enumerate}

\subsection{Results}

\begin{table}[h]
\centering
\caption{Benchmark results (median of 10 runs, in milliseconds)}
\label{tab:benchmarks}
\begin{tabular}{lrrr}
\toprule
Benchmark & Python & Julia & Sounio \\
\midrule
GUM Chain (100K) & 1,244 & 82 & \textbf{41} \\
Monte Carlo (100K) & 871 & 31 & \textbf{19} \\
Cockcroft-Gault (10K) & 4,522 & 210 & \textbf{129} \\
Matrix (50$\times$50) & 892 & 46 & \textbf{23} \\
Transcendental (10K) & 325 & 19 & \textbf{9} \\
\bottomrule
\end{tabular}
\end{table}

\paragraph{Analysis} Sounio achieves consistent 2$\times$ speedup over Julia
due to:
\begin{itemize}
  \item No wrapper object allocation (stack-allocated variance)
  \item SIMD parallelization of value/variance computation
  \item Compile-time resolution of uncertainty propagation rules
\end{itemize}

The 30$\times$ speedup over Python reflects the fundamental overhead of
dynamic typing and operator overloading.

\subsection{Memory Usage}

\begin{table}[h]
\centering
\caption{Peak memory for GUM Chain (100K operations)}
\begin{tabular}{lr}
\toprule
Implementation & Peak Memory \\
\midrule
Python & 847 MB \\
Julia & 124 MB \\
Sounio & \textbf{89 MB} \\
\bottomrule
\end{tabular}
\end{table}

\section{Applications}
\label{sec:applications}

\subsection{Pharmacokinetics}

The Cockcroft-Gault equation for creatinine clearance:
\[
  \text{CrCl} = \frac{(140 - \text{age}) \times \text{weight} \times [0.85 \text{ if female}]}{72 \times \text{SCr}}
\]

In clinical practice, each input has measurement uncertainty. Sounio
automatically propagates these to the output:

\begin{lstlisting}[caption={GUM-compliant drug dosing}]
fn creatinine_clearance(
    age: Knowledge<years>,
    weight: Knowledge<kg>,
    scr: Knowledge<mg_per_dL>,
    is_female: bool
) -> Knowledge<mL_per_min> {
    let factor = if is_female { 0.85 } else { 1.0 }
    let numerator = (Knowledge::exact(140.0) - age) * weight
    numerator.scale(factor) / (Knowledge::exact(72.0) * scr)
}
\end{lstlisting}

The result includes propagated uncertainty suitable for FDA-compliant
documentation.

\subsection{Climate Modeling}

CMIP6 climate projections involve ensemble uncertainty across models.
Sounio's \texttt{ensemble\_mean} correctly computes between-model variance:

\begin{lstlisting}[caption={Climate ensemble analysis}]
fn sea_level_rise_projection(models: &[Knowledge<meters>]) -> Knowledge<meters> {
    Knowledge::ensemble_mean(models)
    // Includes: within-model variance + between-model variance
}
\end{lstlisting}

\subsection{Neuroimaging}

fMRI statistical analysis requires careful propagation of scanner noise
through preprocessing pipelines. Sounio tracks uncertainty through motion
correction, smoothing, and GLM estimation.

\section{Related Work}
\label{sec:related}

\paragraph{Probabilistic Programming}
Stan \citep{carpenter2017stan}, Pyro \citep{bingham2019pyro}, and Gen
\citep{cusumano2019gen} model uncertainty probabilistically but are DSLs
for inference, not systems programming.

\paragraph{Dependent Types}
Languages like Idris \citep{brady2013idris} and F* \citep{swamy2016dependent}
support refinement types but do not model uncertainty semantically.

\paragraph{Units of Measure}
F\# \citep{kennedy2009units} and Sounio both track units at compile time.
Sounio extends this to uncertainty.

\paragraph{Automatic Differentiation}
JAX \citep{bradbury2018jax} and PyTorch \citep{paszke2019pytorch} provide AD
for gradients; Sounio uses similar techniques for uncertainty propagation.

\section{Conclusion}
\label{sec:conclusion}

Sounio demonstrates that uncertainty quantification can be integrated into a
systems programming language with minimal overhead and strong static
guarantees. The \texttt{Knowledge<T>} type makes GUM-compliant uncertainty
propagation as natural as arithmetic, while effect tracking prevents
accidental uncertainty loss.

Future work includes:
\begin{itemize}
  \item Higher-order uncertainty (GUM Supplement 1 Monte Carlo)
  \item Correlated inputs with covariance matrices
  \item GPU-accelerated epistemic computations
  \item Formal verification of propagation correctness
\end{itemize}

\paragraph{Availability}
Sounio is open source under the MIT license.
\begin{itemize}
  \item Repository: \url{https://github.com/Sounio-lang/sounio}
  \item Documentation: \url{https://souniolang.org}
  \item DOI: \href{https://doi.org/10.5281/zenodo.18404188}{10.5281/zenodo.18404188}
\end{itemize}

\bibliographystyle{plainnat}
\bibliography{references}

\end{document}
